\documentclass[../../../include/open-logic-section]{subfiles}

\begin{document}

\olfileid{sth}{choice}{hartogs}
\olsection{القابلية للمقارنة ولمّة هارتوغس}

ذلك هو الجانب الإيجابي. أما الجانب السلبي فهو أن الأمور تصبح
\emph{فوضوية} من دون الاختيار. ولرؤية السبب، إليك نتيجة جميلة
لـ\cite{Hartogs1915}:

\begin{lem}[\emph{في $\ZF$}]\ollabel{HartogsLemma}
لكل مجموعة~$A$، يوجد عدد ترتيبي~$\alpha$ بحيث
$\cardnless{\alpha}{A}$.
\end{lem}

\begin{proof}
إذا كان~$B \subseteq A$ و$R \subseteq B^2$، فإن
$\tuple{B, R} \subseteq V_{\setrank{A}+4}$ بموجب
\olref[ord-arithmetic][using-addition]{rankcomputation}. ومن ثم انظر،
باستعمال الفصل، في:
\[
	C = \Setabs{\tuple{B, R} \in V_{\setrank{A}+5}}{B\subseteq A
	\text{ و$\tuple{B, R}$ ترتيب جيد}}
\]
وباستخدام الاستبدال و
\olref[ordinals][ordtype]{thmOrdinalRepresentation}، كوّن المجموعة:
\[
	\alpha = \Setabs{\ordtype{B, R}}{\tuple{B, R} \in C}.
\]
بموجب~\olref[ordinals][basic]{corordtransitiveord}، تكون~$\alpha$
عددًا ترتيبيًا، لأنها مجموعة متعدية من الأعداد الترتيبية. ففي
النهاية، إذا كان~$\gamma \in \beta \in \alpha$، فإن
$\beta = \ordtype{B, R}$ لبعض~$B \subseteq A$، وعندئذ
$\gamma = \ordtype{B_b, R_b}$ لبعض~$b \in B$ بموجب
\olref[ordinals][iso]{wellordinitialsegment}، ومن ثم
$\gamma \in \alpha$.

للوصول إلى تناقض، افترض وجود !!a{injection}
$f \colon \alpha \to A$. عندئذ، حيث:
\begin{align*}
	B &= \ran{f}\\
	R &= \Setabs{\tuple{f(\alpha), f(\beta)} \in A \times A}{\alpha \in \beta}.
\end{align*}
من الواضح أن~$\alpha = \ordtype{B, R}$ وأن
$\tuple{B, R} \in C$. ومن ثم~$\alpha \in \alpha$، وهذا تناقض.
\end{proof}

يستلزم ذلك نتيجة عميقة:

\begin{thm}[\emph{في $\ZF$}]
الادعاءان الآتيان متكافئان:
\begin{enumerate}
	\item\ollabel{equivwo} مسلّمة الترتيب الجيد؛
	\item\ollabel{equivcompare} إما~$\cardle{A}{B}$ أو
	$\cardle{B}{A}$، لأي مجموعتين~$A$ و$B$.
\end{enumerate}
\end{thm}

\begin{proof}
\emph{\olref{equivwo} $\Rightarrow$ \olref{equivcompare}.} ثبّت~$A$
و$B$. باستدعاء~\olref{equivwo}، يوجد ترتيبان جيدان
$\tuple{A, R}$ و$\tuple{B, S}$. وباستدعاء
\olref[ordinals][ordtype]{thmOrdinalRepresentation}، لتكن
$f \colon \alpha \to \tuple{A, R}$ و
$g \colon \beta \to \tuple{B, S}$ تماثلين. بموجب
\olref[sth][ordinals][basic]{ordinalsaresubsets}، إما
$\alpha \subseteq \beta$ أو~$\beta \subseteq \alpha$. فإذا كان
$\alpha \subseteq \beta$، كانت
$\comp{f^{-1}}{g} \colon A \to B$ !!a{injection}، ومن ثم
$\cardle{A}{B}$؛ وبالمثل، إذا كان~$\beta \subseteq \alpha$، كان
$\cardle{B}{A}$.

\emph{\olref{equivcompare} $\Rightarrow$ \olref{equivwo}.} ثبّت~$A$؛
وبموجب~\olref{HartogsLemma}، يوجد عدد ترتيبي~$\beta$ بحيث
$\cardnless{\beta}{A}$. وباستدعاء~\olref{equivcompare}، يكون
$\cardle{A}{\beta}$. ومن ثم توجد !!{injection}
$f \colon A \to \beta$، ويمكننا استخدام هذا الحقن لترتيب عناصر~$A$
ترتيبًا جيدًا، بأن نعرّف ترتيبًا
$\Setabs{\tuple{a, b} \in A \times A}{f(a) \in f(b)}$.
\end{proof}
\noindent
ونتيجة مباشرة لذلك: إذا فشل الترتيب الجيد، كانت بعض المجموعات
\emph{غير قابلة للمقارنة حرفيًا} من حيث حجمها. فإذا فشل الترتيب الجيد،
كان الحساب الكاردينالي المتناهي التجاوز فوضويًا. فعلينا مثلًا التخلي
عن فكرة أنه إذا كانت~$A$ و$B$ غير منتهيتين، فإن
$\cardeq{A \disjointsum B}{M}$ و$\cardeq{A \times B}{M}$، حيث~$M$ هي
الكبرى بين~$A$ و$B$ (انظر
\olref[card-arithmetic][simp]{cardplustimesmax}). والمشكلة بسيطة: إذا
لم نستطع \emph{مقارنة} حجم~$A$ وحجم~$B$، فلا معنى للسؤال أيهما أكبر.

\end{document}
